\section{The Euler spectral cover and equivariant Hodge atoms}\label{sec:hodgeatoms}

At a generic point of a fixed semisimple germ one can index the distinct Euler values by orbits of the group on primitive idempotents. The definition of a local Hodge atom involves continuing the values over a larger base. The simplest example already distinguishes the two operations. For $X=\mathbb P^1$ with trivial group, multiplication in the hyperplane class satisfies $h\star h=q$. At zero unit coordinate the Euler element is $2h$. At any specified $q\ne0$ there are two eigenvalues, but as $q$ varies their cover is the single connected double cover $s^2=q$. Thus there is one local atom of multiplicity two. We give the general relation between the two labellings.

Use the nonarchimedean parameter field $\mathbb K$ of \cite[\S5.4]{KKPY}, with algebraically closed constant field $\overline{\mathbb Q}$ and positive-valuation parameter $\mathbf y$. On the maximal quantum base of \cite[\S3.5.2]{KKPY} the divisor coordinates are multiplicative characters $q$ with ample valuation. The unit coordinate is unrestricted and the remaining positive-degree parameters are topologically small. Following Cavenaghi, Katzarkov and Kontsevich \cite[Definition~2.3 and \S2.5]{CKK}, restrict to the geometric fixed locus of $G$ and the Hodge group (the kernel of the Tate character). Toric cohomology is generated by algebraic divisor classes and has no odd part. Thus the Hodge group acts trivially and the fixed base is $B_X^G$. Here fixed locus means the equalizer of the action maps. For example, if the divisor torus has character lattice $\Lambda$, its fixed subgroup has character group $\Lambda/\sum_{g\in G}(g-1)\Lambda$, which includes any torsion in the coinvariants.

Let $A$ denote the finite locally free algebra of primary quantum multiplication on this fixed base and let $a$ be the Euler element in the convention of \cite[\S2.1]{CKK}. Write $U\subset B_X^G$ for the locus where multiplication by $a$ has the maximum number of distinct geometric eigenvalues. The local equivariant Hodge atoms are the connected components of the reduced eigenvalue cover $\widetilde U\to U$. Their multiplicities are the covering degrees, taken on each connected base component.

\begin{theorem}\label{thm:atomcover}
Let $X$ be a smooth projective complex toric Fano variety with a finite group $G$ of fan symmetries. The maximum number of distinct Euler eigenvalues on $B_X^G$ is $\dim H^*(X,\mathbb Q)^G$. The full quantum algebra is finite étale over $U$. Its primitive characters form a cover $P=\operatorname{Spec}_U A$, and evaluation of $a$ induces an isomorphism of finite étale covers
\[
                          P/G\ \simeq\ \widetilde U .
\]
Over a connected component of $U$, let $S$ be a geometric fibre of $P$. The local equivariant Hodge atoms are indexed by the étale monodromy orbits on $S/G$. The multiplicity of each atom is the length of that monodromy orbit.
\end{theorem}

\subsection{Separation on the fixed base}
We first show that the bound in the theorem is attained. The distinction between two Euler conventions matters away from the small divisor slice. In primary coordinates the element used in Proposition~\ref{prop:cluster} is $E_+=c_1(X)+\sum_p(1-p)t_{2p}$, where $t_{2p}$ denotes the component in $H^{2p}$. The expression printed in \cite[\S2.1]{CKK} and \cite[(3.20), (3.31)]{KKPY} is $E_-=c_1(X)+\sum_p(p-1)t_{2p}$. A simultaneous negation of an entire multiplication operator merely changes the spectral coordinate $\lambda$ to $-\lambda$. The relative sign of the nondivisor terms in these two expressions calls for a separate check. We use only their multiplication in the primary algebra.

On a simply connected semisimple chart let $u_s$ be the canonical coordinates for $E_+$. The coefficient of any tangent vector $v$ on the primitive factor $s$ is $du_s(v)$. Proposition~\ref{thm:main} supplies a semisimple point on the complex fixed locus on the large-radius logarithmic branch. The independent orbit coordinates from Section~\ref{sec:objects} then separate distinct orbits generically for $E_+$. They do so for $E_-$ as well. Indeed, for two labels in different orbits put $f=u_s-u_t$ and let $L$ be the vector field with component $(1-p)t_{2p}$ in degree $2p$, $p\ge2$, and zero in the other degrees. The unit derivative of $f$ vanishes. The difference of the $E_-$ eigenvalues at those two labels is therefore $f-2L(f)$.

Suppose this difference vanished identically on the fixed germ. Choose a starting point with $f=c\ne0$. At the same small divisor parameter, scale $t_{2p}$ by $b^{p-1}$ for $p\ge2$. The convergence polydisc may be chosen about zero in these variables, so the arc extends through $b=0$ and a neighbourhood of $b=1$. The scaling stays on the fixed locus. Near $b=1$ we would have $f+2b f'=0$, giving $f=c/\sqrt b$. Let $M$ denote multiplication by $E_+$ on the full cohomology. The characteristic polynomial of $M\otimes 1-1\otimes M$ is monic, holomorphic along the whole arc, and has $f$ as a root near $1$. Substituting $b=v^2$, the supposed root $c/v$ would then solve this polynomial on a punctured disk about zero by the identity theorem. The leading pole from its monic term makes this impossible. This argument requires semisimplicity only near the chosen starting point. Avoiding the finitely many proper equalities for both elements gives a common separated open. On it every orbit of primitive factors accounts for one eigenvalue for either convention.

Put $r=\dim H^*(X,\mathbb Q)^G$. At these complex points $r$ is the number of orbits, since the primitive idempotents are permuted. To attain the same number over $\mathbb K$, work in rational coordinates $D$ on the invariant divisor space and $\eta$ on the invariant sum of the $H^{2p}$ for $p\ge2$. Put the unit coordinate equal to zero, which leaves all eigenvalue differences unchanged. Make the specialization
\[
         q^\beta=x^{c_1(X)\cdot\beta}\exp\langle\beta,D\rangle,
                     \qquad t_{\ge4}=\eta .
\]
The primary structure constants are rational series in $x,D,\eta$. This follows from the divisor equation, rationality of the invariants and finiteness of effective classes in bounded anticanonical degree. Section~\ref{sec:exist} gives their complex convergence near the origin. The specialization with $x\ne0$ covers a neighbourhood in the fixed branch at its large-radius anticanonical points, since $D$ supplies all local fixed divisor directions there.

In this rational series ring take the determinant of the trace form of the full algebra, multiplied by the discriminants of the two Euler operators on the $r$-dimensional invariant algebra. This product is nonzero by the common complex separated open. For suitable rational constant vectors $D_1,\eta_1$ the substitution $D=xD_1$, $\eta=x\eta_1$ leaves it nonzero: its lowest nonzero homogeneous part evaluated at $(1,D_1,\eta_1)$ is a nonzero polynomial in those vectors. Evaluate the resulting formal arc at $x=\mathbf y$ in $\mathbb K$. All series are defined, $q$ has ample valuation, and the full algebra is semisimple with $r$ distinct invariant eigenvalues for each convention. The case with no nondivisor directions is included.

This arc also compares the geometric sets of labels. Over the fraction field of the one-variable series an element with rational constant coordinates in a cohomology basis separates the primitive characters: avoid the finite union of proper linear conditions imposed by their pairwise differences. Its multiplication polynomial is monic with rational convergent coefficients. After a finite substitution $x=w^m$ the roots have convergent power series in $w$, by continuation of the simple roots on the punctured disk and boundedness at zero. Their expansions agree with the formal Newton--Puiseux roots in $\overline{\mathbb Q}[[w]]$. Interpolation in this element gives the idempotents with meromorphic coefficients in $w$. The same constant rational matrices act by $G$ on these expressions in either field. At a small positive complex value and the corresponding Puiseux branch the primitive sets are therefore isomorphic as $G$-sets. Within the connected complex parameter polydisc with $x\ne0$ one can join the two relevant complex semisimple points while avoiding the proper analytic discriminant. Such a path compares with the primitive labels of Proposition~\ref{thm:main} on the chosen branch. The construction makes no choice of labels on a different fixed torus component.

\subsection{Semisimplicity over the maximum locus}
\begin{proof}[Proof of Theorem~\ref{thm:atomcover}]
In a geometric fibre the minimal polynomial of $a$ on $A$ is its minimal polynomial on $A^G$ \cite[Lemma~5.19]{KKPY}. In either algebra the operator $p(a\star)$ vanishes exactly when its value on the unit is zero. Thus $r$ bounds the number of eigenvalues everywhere, also on components outside the logarithmic chart. The preceding arc attains this bound. At every point of $U$ the $r$ powers $1,a,\ldots,a^{r-1}$ form a basis of $A^G$ and the minimal polynomial is separable. It follows that $A^G$ is finite étale on $U$.

Here the Frobenius pairing allows one to pass from the invariants to the whole algebra. Over an algebraically closed residue field, decompose the fibre as a product of local Artin algebras and write $\ell(v)=(v,1)$. The functional $\ell$ is invariant under $G$ and its pairing restricts nondegenerately to each local factor. In such a factor $C$ with maximal ideal $\mathfrak m$ the pairing identifies $\mathrm{Ann}(\mathfrak m)$ with the dual of $C/\mathfrak m$. Its restriction under $\ell$ is consequently an isomorphism from this one-dimensional socle to the field. If $C$ is not a field, the socle lies in $\mathfrak m$. Its vector normalized by $\ell(v)=1$ is nilpotent and fixed by the stabilizer of $C$. The sum of its translates, one in each distinct orbit factor, is a nonzero invariant nilpotent. This contradicts reducedness of the invariant fibre, since averaging commutes with scalar extension. Thus every geometric fibre of $A$ is étale. The algebra is finite locally free, so $P$ is finite étale.

Étale locally write $A$ as the functions on $S$. Then $A^G$ is the functions on $S/G$, and $a$ takes distinct values on distinct orbits. The characteristic polynomial on $A$ has one factor for each orbit, to the power of its cardinality. Reduction leaves one copy, identifying its zero locus canonically with $\operatorname{Spec}_U A^G=P/G$. The deck action of $G$ commutes with monodromy because it acts over the base. The connected components of a finite étale cover correspond to the monodromy orbits on a geometric fibre. This proves the formula for their degrees.
\end{proof}

For instance, in the projective-line example the divisor parameter ranges over $0<|q|<1$ and the unit parameter $t_0$ translates both eigenvalues by the same scalar. The whole reduced cover has coordinates $(s,t_0)$ with $0<|s|<1$, and projects by $(s,t_0)\mapsto(s^2,t_0)$. Its source is connected, as is seen by exhausting the punctured disk times the affine line by products of closed annuli and disks with integral affinoid algebras. On the anticanonical slice $q=x^2$ the cover splits into the two sections $s=\pm x$. One can thus have separate labels in the logarithmic family even when they form one atom on the multiplicative base.

\subsection{A family with trivial monodromy}\label{sec:centralatoms}
For the centrally symmetric del Pezzo varieties of Section~\ref{sec:example} the entire primitive cover can be split. In dimension four, for instance, the two orbits of cardinality two have isomorphic stabilizers and could in principle be exchanged by monodromy. To distinguish their sheets we use the grading of the quantum product.

\begin{proposition}\label{prop:centralatoms}
For the centrally symmetric del Pezzo variety of dimension $2d$, $d\ge1$, let the full fan group act. The primitive cover is constant as a cover with group action on the entire nonarchimedean fixed base. There are $(d+1)(d+2)/2$ local equivariant Hodge atoms, each of covering degree one. Write $m=2d+1$. For each $d+1\le k\le m$ there is one atom with eigenspace dimension $\binom{m}{k}$ and there are $(2k-m-1)/2$ atoms with eigenspace dimension $2\binom{m}{k}$, counted with these labels if dimensions coincide.
\end{proposition}
\begin{proof}
The character lattice of the divisor torus consists of the integral relations among the rays. Write a relation on the opposite pairs as $(b_i+c,b_i)$, since the only relation among the $v_i$ is their sum. Permuting the pairs gives coinvariant coordinates $(B,c)$ for $B=\sum_i b_i$. Inversion sends these to $(B+mc,-c)$. Thus the full coinvariant lattice is $\Z^2/\Z(m,-2)$. Since $m$ is odd, the anticanonical degree $2B+mc$ identifies this lattice with $\Z$. The entire fixed divisor torus has a unique coordinate $x$ with $q^\beta=x^{c_1\cdot\beta}$. Its valuation domain is $0<|x|<1$. The fixed base, including the invariant higher inputs and the unit line, is connected.

Let $e_i$ denote a homogeneous basis vector of degree $2i$, and use the global frame $e'_i=x^{-i}e_i$, suppressing multiplicity indices. Write $\eta_p$ for the parameters of degree $2p$ with $p\ge2$. A contribution of degree $\beta$ to the $e_k$ coefficient of $e_i\star e_j$, with further inputs of half-degrees $p_1,\ldots,p_l$, satisfies the dimension equation
\[
             c_1\cdot\beta=i+j-k+\sum_{a=1}^l(p_a-1).
\]
Hence in the primed frame the structure constants are power series in $u_p=x^{p-1}\eta_p$ alone. Their coefficients are rational. For prescribed input multiplicities the anticanonical degree is fixed, admitting only finitely many effective curve classes. These formal series define analytic functions on the open unit polydisc since the constant field is trivially valued. All allowed fixed parameters have $|u_p|<1$. The unit coordinate drops from the product. Subtracting its scalar contribution from the Euler element and dividing by $x$ likewise gives an element over this power series ring, with constant term $c_1$ in the normalized small algebra.

At $u=0$ the structure constants in this frame are those of the small product with each Novikov monomial set to one. Its toric presentation sends the divisor of each ray to the corresponding Laurent monomial, hence their sum $c_1$ to $\sum_i(y_i+y_i^{-1})$ \cite[Proposition~3.3]{OT}, \cite[\S1.7]{MM}. This identification respects the fan permutations. The monotone small algebra for this entire family was proved semisimple in \cite[Proposition~1.6.1]{MM}. The equations in Section~\ref{sec:example} give its primitive characters over $\overline{\mathbb Q}$. Each special idempotent lifts uniquely in the finite free algebra over the complete power series ring on the invariant $u$-coordinates. Indeed the linearization of $e^2=e$ at its reduction is multiplication by $2e-1$, an involution. The lifts are orthogonal and sum to one by uniqueness; they form a basis by reduction. Their coefficients, being constants in $\overline{\mathbb Q}$ at each power of $u$, converge on the same open polydisc. Pulling back to the full fixed base gives a simultaneous splitting with the original permutation action.

We check that distinct orbits have different normalized Euler values throughout this base. At $u=0$ the orbit with a majority of $k$ coordinates equal to $r$ has value $j(r+r^{-1})$, where $j=2k-m$ is positive and odd and $r^j=(-1)^{m-k}$. For fixed $j$, equality of two values means that their roots form the same inverse pair. If $r$ is nonreal, $r^2$ has odd order $a>1$ and $r+r^{-1}=r^{-1}(1+r^2)$ is an algebraic unit. In fact squaring permutes the primitive $a$th roots, so cancelling the products of $1-z$ against the products of $1-z^2$ gives $\prod_z(1+z)=1$. For real $r$ the factor is $\pm2$. Equality for two nonreal roots with different $j$ would make the positive rational $j/j'$ an algebraic unit. Two real roots of equal value also force $j=j'$. In the mixed case the same unit argument would make one odd integer twice another. Thus the constant terms separate all distinct orbits. Differences with nonzero constant term remain nonzero at every allowed $u$ by the ultrametric inequality.

There is one real root for each $j$ and there are $(j-1)/2$ nonreal inverse pairs. The group permutes the majority set transitively and simultaneous inversion exchanges the two roots of each nonreal pair. This gives the stated eigenspace dimensions. Summing $(j+1)/2$ over the positive odd numbers $j\le m$ gives the number of separate sheets.
\end{proof}

For the degree-six surface this gives three atoms with eigenspace dimensions $1,2,3$. For the fourfold there are six, with dimensions $1,2,2,5,10,10$. Even the two isomorphic two-element orbits lie on separate components. The ordinary exceptional collection in this family and its fan action are known \cite[Proposition~4.12]{BDM}; see also Castravet--Tevelev \cite[Theorem~6.6]{CTI}. The relation between the three surface blocks and the small spectrum is discussed in \cite[Remark~1.6]{ESS}. Proposition~\ref{prop:centralatoms} computes the degrees on the full fixed quantum base.


\subsection{Categories on the orbit cover}
For a component $C$ of $\widetilde U=P/G$, let $p_C\colon P_C=P\times_{P/G}C\to C$ be its inverse image. The finite étale algebra $(p_C)_*\mathcal O_{P_C}$ has a canonical $G$-action. Its perfect modules form a family whose fibre at a geometric point is the category of perfect complexes on a transitive $G$-set $O$. Monodromy can exchange only isomorphic transitive $G$-sets. For an atom of covering degree $d$ with $|O|=m$, a geometric base fibre has $d$ eigenvalues in that component, each with eigenspace dimension $m$ and invariant dimension one.

Whenever a generic complex fixed parameter carries a marked exceptional collection as in Proposition~\ref{prop:blocktransport}, Proposition~\ref{prop:stabilizers} identifies its individual orbit blocks, with their $G$-actions, with the complex models of these same finite sets after a choice of the label comparison above on the anticanonical component. This describes their abstract equivariant categories. To identify embedded subcategories of $D^b\mathrm{Coh}(X)$ under continuation one must also compare the phase-ordered mutations. Compatibility with the global equivalence of Hodge atoms further involves blowups and projective bundles as in \cite[\S2.5]{CKK}.
